\section{Conclusões}

Este trabalho apresentou um protótipo Web voltado à recomendação para a formação de equipes em projetos de desenvolvimento de software. Ao longo do texto, evidenciou-se que a alocação de profissionais, quando realizada de forma manual e baseada em percepções subjetivas, pode resultar em times desbalanceados e comprometer os prazos e a qualidade das entregas organizacionais. Diante desse cenário, a pesquisa estruturou um sistema que integra algoritmos de aprendizado de máquina, responsáveis pelo cálculo quantitativo de afinidade, com Grandes Modelos de Linguagem (LLMs), encarregados de processar interações e justificar as escolhas geradas de forma qualitativa.

O desenvolvimento contemplou a fundamentação teórica e a análise de trabalhos correlatos, que estabeleceram o referencial necessário para a compreensão de sistemas de recomendação, reconhecimento de padrões e interfaces conversacionais. A especificação do sistema foi feita por meio da modelagem estrutural e funcional, com uma arquitetura baseada em microsserviços e uma abordagem híbrida de entrada de dados, na qual o gestor pode utilizar formulários ou conversar com o sistema em linguagem natural. A estrutura relacional de dados foi modelada para o mapeamento quantitativo e qualitativo de competências, perfis de desenvolvedores e requisitos de projetos, com um catálogo de habilidades que corresponde exatamente aos atributos reconhecidos pelo modelo preditivo.

O motor preditivo foi treinado em Python com a biblioteca PyCaret, a partir dos dados da Stack Overflow Developer Survey 2024, complementados com dados sintéticos de competências comportamentais e de disponibilidade. Entre os 14 algoritmos comparados pela função \texttt{compare\_models}, o LightGBM apresentou o melhor desempenho, com $R^2$ de 0,9333 e MAE de 3,33 pontos na validação cruzada, resultado próximo do limite imposto pelo ruído da própria base. O modelo ajustado foi salvo como um arquivo e passou a ser carregado por um serviço em FastAPI, consultado pelo \textit{backend} a cada recomendação. A integração com o modelo de linguagem foi feita com o Spring AI e a API do Gemini, que atua em duas pontas. Na entrada, converte o texto livre do gestor em parâmetros de busca e, na saída, redige a justificativa da equipe a partir do ranqueamento completo dos candidatos, o que permite explicar também por que um perfil ficou de fora.

A avaliação foi conduzida em quatro cenários controlados, nos quais as equipes recomendadas foram comparadas com 100 equipes sorteadas entre os mesmos candidatos elegíveis. Nos cenários em que havia escolha a ser feita, o protótipo superou a composição aleatória em todas as métricas avaliadas. A cobertura de papéis chegou a 100\%, contra 59,60\% e 78,33\% da linha de base nos Cenários 1 e 2, e o escore médio da equipe ficou cerca de 3,3 desvios padrão acima da média dos sorteios no Cenário 2. Em contexto de escassez, o sistema entregou uma equipe parcial acompanhada da descrição do que faltou, em vez de recusar a solicitação. Com dois projetos disputando os mesmos profissionais, o desconto das horas já confirmadas eliminou a sobreposição entre as equipes sem impedir a formação completa da segunda.

Em síntese, os resultados indicam que a integração entre o motor preditivo e as regras de composição produz recomendações consistentes com os requisitos declarados e superiores a alocações não orientadas por dados. A principal contribuição do trabalho está na combinação entre o ranqueamento quantitativo, feito pelo modelo preditivo, e a explicação em linguagem natural, que torna esse ranqueamento compreensível para o gestor. Essa combinação responde ao problema apontado na introdução, segundo o qual recomendações geradas puramente por Inteligência Artificial são difíceis de interpretar e diminuem a confiança no sistema. O protótipo se converte, assim, em um assistente técnico orientado pela transparência, no qual a decisão final continua com o gestor, que recebe uma sugestão justificada e decide se confirma a equipe.

Os resultados, contudo, se restringem à validação técnica da solução. O escore de afinidade foi calculado de forma sintética, devido à indisponibilidade de bases públicas com o histórico real de sucesso de equipes, e os perfis utilizados nos cenários são fictícios. Além disso, a composição aleatória demonstra que o desempenho do protótipo não decorre do acaso, mas não equivale ao desempenho de um gestor experiente.

Como trabalhos futuros, destaca-se a aplicação do protótipo em ambiente organizacional real, com o registro histórico das alocações e de seus resultados. Esses dados permitiriam substituir a variável alvo sintética por dados observados, treinar novamente o modelo e comparar as equipes recomendadas com as formadas por gestores. Outro caminho é complementar a métrica de dispersão com um critério de diversidade de experiência, já que no Cenário 2 a equipe inteiramente sênior obteve o melhor valor possível, o que não atende projetos nos quais a formação de profissionais menos experientes também seja um objetivo. Espera-se que esses desdobramentos aproximem a solução do uso cotidiano e reforcem sua viabilidade para atenuar as incertezas atreladas à gestão e à alocação de equipes de desenvolvimento de software.
